\chapter{Variables, Distributions \& the Mapper}
\label{ch:variables}

\newthought{Variables define the space you scan.} Each entry under
\yamlkey{Sampling.Variables} names one sampling coordinate, and its
\yamlkey{distribution} fixes how a u-space coordinate in $[0,1]$ becomes a physical
value. Sometimes that is the whole story. Sometimes the coordinates you want to
\emph{sample} are not the parameters you want to \emph{use} --- and that is what
\yamlkey{Sampling.Mapper} is for.

\section{Two ways to describe a space}
\label{sec:mapper-modes}

There are exactly two shapes a \yamlkey{Sampling} block can take here, and picking
between them is the first decision of this chapter.

\begin{description}[style=nextline, leftmargin=1.2em]
  \item[\textnormal{\textbf{Variables only.}}] Every parameter your workflow needs is
        something you can sample directly. One \yamlkey{Variables} entry per parameter,
        each with its own distribution, and nothing else. This is the common case and
        the whole of Sections~\ref{sec:variables-entry} and~\ref{sec:distributions}.
  \item[\textnormal{\textbf{Variables + Mapper.}}] You sample in one set of coordinates
        and hand the workflow another. \yamlkey{Variables} declares what the sampler
        moves through; \yamlkey{Mapper} derives the rest with expressions
        (Section~\ref{sec:mapper}).
\end{description}

The second mode earns its keep when the parameters you care about are
\emph{constrained}. Suppose two masses must sum to a fixed scale. Declaring \code{m\_1}
and \code{m\_2} as independent variables and then rejecting everything off the constraint
with \yamlkey{selection} wastes almost the entire box. Sampling the sum and a ratio
instead puts every point on the surface by construction:

\begin{yamlwin}
Sampling:
  Method: Random
  Variables:
    - name: m_sum
      distribution: {type: Flat, parameters: {min: 200.0, max: 2000.0}}
    - name: r
      distribution: {type: Flat, parameters: {min: 0.0, max: 1.0}}
  Mapper:
    - name: m_1
      expression: "m_sum * r"
    - name: m_2
      expression: "m_sum * (1 - r)"
\end{yamlwin}

The scan is still two-dimensional --- \yamlkey{Mapper} adds parameters, never
dimensions --- but \code{m\_1} and \code{m\_2} arrive at the workflow as ordinary
parameters, and no point is ever thrown away.

\begin{jtip}
The question to ask is not ``do I want derived quantities?'' but ``\emph{where does the
sampler need to be uniform?}'' If a derived quantity is only ever \emph{read} --- a ratio
you want plotted, a combination entering the likelihood --- an \yamlkey{Operas} module or
a likelihood term does the job just as well. Reach for \yamlkey{Mapper} when the
\emph{geometry} of the scan is the point.
\end{jtip}

\section{The \yamlkey{Variables[]} entry}
\label{sec:variables-entry}

\begin{yamlwin}
Sampling:
  Variables:
    - name: m0                     # required
      description: "unified scalar mass"   # informational only
      distribution:
        type: Log                  # default: Flat
        parameters: {min: 100.0, max: 10000.0}
\end{yamlwin}

\begin{keylist}
  \item[name] \req\ in practice --- an entry with an empty or missing name is
        \emph{silently dropped}. The name becomes the observable key and the expression
        symbol for this parameter everywhere downstream.
  \item[description] \opt, purely informational.
  \item[distribution.type] \opt, default \code{Flat}. One of the five types below.
  \item[distribution.parameters] the per-type parameter map --- and the home of the
        method-specific extras \yamlkey{num} (\sampler{Grid}) and \yamlkey{length}
        (\sampler{Bridson}).
\end{keylist}

\begin{jwarn}
A missing required distribution parameter surfaces at sampler startup as a raw
\code{KeyError} --- not as a schema message naming the variable. If a scan dies
instantly with \code{KeyError: 'max'}, check your \yamlkey{parameters} maps first.
\end{jwarn}

\section{The distribution catalog}
\label{sec:distributions}

\begin{table}[ht]
  \begingroup
  \setlength{\tabcolsep}{0pt}
  \begin{tabular}{@{}p{0.18\linewidth}p{0.26\linewidth}p{0.56\linewidth}@{}}
    \toprule
    \textbf{type} & \textbf{Parameters} & \textbf{Mapping $u\in[0,1] \to x$} \\
    \midrule
    \code{Flat}       & \yamlkey{min}, \yamlkey{max} & $x = \yamlkey{min} + (\yamlkey{max}-\yamlkey{min})\,u$ \\
    \code{Log}        & \yamlkey{min}, \yamlkey{max} ($>0$) & $x = \exp\!\big(\ln \yamlkey{min}$ \allowbreak $+ (\ln \yamlkey{max} - \ln \yamlkey{min})\,u\big)$ \\
    \code{Normal}     & \yamlkey{mean}, \yamlkey{stddev} & $x = \yamlkey{mean} + \yamlkey{stddev}\cdot\Phi^{-1}(u)$ \\
    \code{Log-Normal} & \yamlkey{mean}, \yamlkey{stddev} & $x = \exp\!\big(\yamlkey{mean}$ \allowbreak $+ \yamlkey{stddev}\cdot\Phi^{-1}(u)\big)$ \\
    \code{Logit}      & \yamlkey{location} (0), \yamlkey{scale} (1) & $x = \mathrm{logit}(u)\cdot\yamlkey{scale} + \yamlkey{location}$ \\
    \bottomrule
  \end{tabular}
  \endgroup
  \caption{Distribution types. $\Phi^{-1}$ is the standard normal quantile;
  $\mathrm{logit}(u)=\ln\frac{u}{1-u}$. \code{Flat}/\code{Log} parameters are required;
  \code{Logit}'s default as shown.}
  \label{tab:distributions}
\end{table}

Intuition for choosing:

\begin{itemize}
  \item \code{Flat} --- the default; equal prior weight per unit of $x$.
  \item \code{Log} --- equal weight per \emph{decade}; the right choice for masses and
        couplings spanning orders of magnitude.
  \item \code{Normal} / \code{Log-Normal} --- concentrate points around a known central
        value, e.g.\ a measured nuisance parameter.
  \item \code{Logit} --- an unbounded variable sampled through a squashing map; heavy
        tails compared to \code{Normal}.
\end{itemize}

\paragraph{Method extras on the same map.} \sampler{Grid} requires \yamlkey{num} (points
along this dimension) and \sampler{Bridson} requires \yamlkey{length} (the u-space box
edge for this dimension) inside \yamlkey{distribution.parameters} --- geometry keys that
ride along with the distribution parameters. Details in
Chapters~\ref{ch:grid} and~\ref{ch:bridson}.

\section{The \yamlkey{Sampling.Mapper} block}
\label{sec:mapper}

\yamlkey{Mapper} is an ordered \emph{list}, one entry per derived parameter, and each
entry has exactly two keys:

\begin{yamlwin}
Sampling:
  Variables:
    - name: t
      distribution: {type: Flat, parameters: {min: 0.0, max: 6.283}}
  Mapper:
    - name: x
      expression: "cos(t)"
    - name: y
      expression: "sin(t)"
\end{yamlwin}

\begin{keylist}
  \item[name] \req. The derived parameter's name --- the symbol it gets everywhere
        downstream.
  \item[expression] \req. How to compute it, in the shared language of
        Chapter~\ref{ch:expressions}, so the full function catalog and every \Operas\
        namespace are available.
\end{keylist}

Those two keys are the whole entry; anything else is rejected. Note the shape this
gives the block: your parameter names are \emph{values} of \yamlkey{name}, never
\textsc{yaml} keys. \yamlkey{Mapper} therefore reads like
\yamlkey{Variables} --- a list of named things --- rather than like a settings map, which
is the right way to think about it, since both lists describe parameters and neither
describes options.

\subsection{What happens, in order}
\label{sec:mapper-order}

For every point the runtime does two steps, always in this order:

\begin{enumerate}[leftmargin=1.6em]
  \item \textbf{Distributions.} Each u-coordinate becomes a physical value through its
        variable's \yamlkey{distribution} (Table~\ref{tab:distributions}). This step
        happens whether or not you wrote a \yamlkey{Mapper}.
  \item \textbf{Derivations.} Each \yamlkey{Mapper} expression is evaluated and its
        result added to the same dictionary, so later entries can use earlier ones.
\end{enumerate}

What comes out is one flat set of parameters --- the sampling variables \emph{and} the
derived names together. From that moment they are indistinguishable: both are seeded into
the observables dictionary, both are visible to \yamlkey{selection}, to calculators and
\yamlkey{Operas} modules, to the likelihood, and both land in the output records.
\JPlot\ even prefers \yamlkey{Mapper} names when it picks default plot axes, on the
assumption that if you bothered to derive them, they are the coordinates you think in.

\begin{jnote}
The control process and every Worker run the \emph{same} mapping implementation over the
same u-coordinates. A point's parameters therefore do not depend on where they were
computed --- which is what lets a resumed scan rebind a \code{uuid} to exactly the
coordinates it had before.
\end{jnote}

\subsection{Order does not matter (dependencies do)}

\yamlkey{Mapper} entries may reference each other, and the runtime sorts out the order
itself by following the dependencies. This list works exactly as written, even though
\code{z} is listed before the two names it uses:

\begin{yamlwin}
Mapper:
  - name: z
    expression: "x + y"    # evaluated last anyway
  - name: x
    expression: "cos(t)"
  - name: y
    expression: "sin(t)"
\end{yamlwin}

What is not allowed is a cycle --- \code{x} defined from \code{y} while \code{y} is
defined from \code{x} --- which is reported at load time as \code{JV2-MAP-004} naming the
entries involved.

\subsection{A closed namespace}
\label{sec:mapper-namespace}

A \yamlkey{Mapper} expression may refer to \textbf{sampling variables and other
\yamlkey{Mapper} names, and nothing else}. Observables produced by calculators or
\yamlkey{Operas} modules are not visible; neither is \code{LogL}, nor \code{uuid}.

This is a real restriction and it is deliberate. The mapping runs \emph{before} anything
has been computed --- it is what produces the parameters those calculators will be given.
Letting it read their output would be circular. Referencing an unknown name is
\code{JV2-MAP-002}, and the message lists every name that \emph{was} available, which is
usually enough to see the mistake.

Derived names must also be new:

\begin{keylist}
  \item[not already taken] A derived name cannot shadow a \yamlkey{Variables} name, and
        two \yamlkey{Mapper} entries cannot share a name (\code{JV2-MAP-003}).
  \item[not a reserved constant] \code{Pi}, \code{pi}, \code{PI}, \code{E}, \code{Inf}
        (\code{JV2-MAP-005}).
  \item[not a reserved column] \code{uuid}, \code{sample\_index}, \code{status},
        \code{LogL} (\code{JV2-MAP-006}).
\end{keylist}

\begin{jwarn}
One name clash is only a \emph{warning}, and it is the one worth watching:
\code{JV2-MAP-052} fires when a \yamlkey{Mapper} name is also produced downstream as an
\yamlkey{Operas} output or a calculator \yamlkey{Dump} variable. The scan runs, but the
downstream value overwrites the observable while the recorded parameter keeps the mapped
one --- the same name carrying two different numbers. Rename one of them.
\end{jwarn}

\subsection{Dimensions, priors, and which methods care}
\label{sec:mapper-priors}

The sampler moves through one dimension per \yamlkey{Variables} entry. \yamlkey{Mapper}
never changes that count, so \sampler{Grid}'s \yamlkey{num} and \sampler{Bridson}'s
\yamlkey{length} still belong to the variables, and a card with two variables and five
\yamlkey{Mapper} entries is still a two-dimensional scan.

For the coverage methods --- \sampler{Random}, \sampler{Grid}, \sampler{Bridson},
\sampler{AdaptiveBridson} --- that is the end of it. For the statistical methods
(\sampler{Dynesty}, \sampler{MultiNest} and the \textsc{mcmc} family) there is one thing
to be careful about, and the validator will say so:

\begin{keylist}
  \item[JV2-MAP-050] More \yamlkey{Mapper} parameters than sampling variables. Your
        points lie on a lower-dimensional surface inside the space you appear to be
        exploring --- fine for a coverage scan, rarely what you want for inference.
  \item[JV2-MAP-051] A nonlinear \yamlkey{Mapper} expression. The prior is defined by
        \yamlkey{Variables}, not by the derived parameters, and a nonlinear map does not
        carry a flat prior across unchanged.
\end{keylist}

Both are warnings, not errors --- the scan runs. If the derived parameters really are
your inference coordinates, the fix is yours to make: include the absolute Jacobian of
the transform as a term in \yamlkey{LogLikelihood} (Chapter~\ref{ch:loglikelihood}).

\begin{jnote}
\sampler{CSV} does not accept a \yamlkey{Mapper} (\code{JV2-MAP-010}). Its points are
already physical --- they come from the file's columns --- so there is no
u\,$\rightarrow$\,x step to reparameterize. Derive what you need as extra columns in the
\textsc{csv} itself.
\end{jnote}

\subsection{Changing a \yamlkey{Mapper} ends a resume}

The checkpoint records a fingerprint of the \yamlkey{Mapper} entries together with the
variable distributions. If either changed, \cli{-{}-resume} is \textbf{refused} and the
stale checkpoint is removed, so a later run cannot quietly pick it up. Reordering the
list is not a change --- the fingerprint is taken over the entries by name, and the
evaluation order comes from the dependencies anyway --- but editing a name or an
expression is.

The reason is worth stating plainly: a resumed scan replays \code{uuid}s against
coordinates. Under an edited mapping the same \code{uuid} would name a different point in
parameter space, and the finished half of your scan would no longer mean what its records
say. Editing a \yamlkey{Mapper} means starting a fresh scan --- or restoring the original
card.

\subsection{Diagnostics at a glance}

\begin{table}[ht]
  \footnotesize
  \begingroup
  \setlength{\tabcolsep}{0pt}
  \begin{tabular}{@{}p{0.2\linewidth}p{0.15\linewidth}p{0.65\linewidth}@{}}
    \toprule
    \textbf{Code} & \textbf{Level} & \textbf{Meaning} \\
    \midrule
    \code{JV2-MAP-001} & error & \yamlkey{Mapper} is not a list of
      \yamlkey{name}/\yamlkey{expression} entries, or an entry is missing one of them \\
    \code{JV2-MAP-002} & error & expression references a name outside the closed
      namespace \\
    \code{JV2-MAP-003} & error & derived name shadows a \yamlkey{Variables} name or
      repeats an earlier \yamlkey{Mapper} name \\
    \code{JV2-MAP-004} & error & dependency cycle among \yamlkey{Mapper} entries \\
    \code{JV2-MAP-005} & error & derived name is a reserved constant \\
    \code{JV2-MAP-006} & error & derived name is a reserved output column \\
    \code{JV2-MAP-007} & error & the expression does not parse \\
    \code{JV2-MAP-010} & error & \yamlkey{Mapper} used with \sampler{CSV} \\
    \code{JV2-MAP-050} & warning & more derived parameters than sampling dimensions,
      under a statistical method \\
    \code{JV2-MAP-051} & warning & nonlinear reparameterization under a statistical
      method \\
    \code{JV2-MAP-052} & warning & derived name also produced downstream \\
    \bottomrule
  \end{tabular}
  \endgroup
  \caption{\yamlkey{Sampling.Mapper} diagnostics. Errors are reported at load time,
  before any Worker starts.}
  \label{tab:map-codes}
\end{table}

\section{Practical guidance}

\begin{jtip}
Declare every variable with the distribution you \emph{mean}, not \code{Flat} plus a
transformation somewhere else. The declared distribution is what samplers use for
geometry (\sampler{Bridson} radii, \sampler{AdaptiveBridson} refinement), what sets the
prior, and what your records show --- keeping the intent in one place. In particular, a
\code{Log} variable is not the same scan as a \code{Flat} variable with an
\code{exp()} in the \yamlkey{Mapper}: the first is uniform per decade, the second is
uniform per unit, and only the first is what you meant when you said ``scan this mass
logarithmically''.
\end{jtip}

Names to avoid, for variables and derived parameters alike: single letters that collide
with mathematical constants and functions in the expression language (\code{E},
\code{I}, \code{N}, \code{pi}, \code{beta}, \code{gamma}); Chapter~\ref{ch:expressions}
lists the reserved names. For \yamlkey{Mapper} entries the hard cases are rejected at
load time (Table~\ref{tab:map-codes}), but a variable named \code{beta} will simply
behave strangely.

\begin{jtip}
Start with \yamlkey{Variables} alone. Add a \yamlkey{Mapper} when you can name the
constraint it enforces --- ``the two masses sum to the scale'', ``the point lies on the
unit circle'' --- and keep the derivations short enough to read. A \yamlkey{Mapper} that
has grown into a calculation belongs in an \yamlkey{Operas} module instead, where it can
be tested.
\end{jtip}
